\begin{claim}{C.55}{Proved}{The full local centralizer and its exact global limits}
Retain the ambient-null based plane loop $f_t$, its generated
group $G$, its commuting operator group $H$, and
$L=\overline G$ from C.42. On every open $G$-orbit define
the closed topological group
\[
 \mathcal C_O=C_{\operatorname{Homeo}(O)}(G|_O),
 \qquad H_O\longrightarrow\mathcal C_O,\quad h\longmapsto h|_O.
                                                        \tag{C.55.1}
\]
The topology is compact-open, including inverse convergence.
The kernel of the displayed homomorphism consists exactly
of operators fixing $O$ pointwise. Its image closure is
the previously used $Q_O$, a closed subgroup of
$\mathcal C_O$.

Let $U^\ddagger$ be the union of open orbits with
$\dim\mathcal C_O=2$ or
$|\mathscr F_O^{\rm all}|\geq3$, using C.53's full global
quotient family, and put $X^\ddagger=\overline{U^\ddagger}$.
There is a jointly continuous based contraction on
$X^\ddagger$, with jointly continuous inverse and
commutation with $H$, extending the specified contraction
on $X^\dagger$ in C.53. Its nonopen-orbit values are the
exponential branch of C.52.1.

Choose finitely many open orbits $O_1,\ldots,O_m$
disjoint from $X^\ddagger$ for which
$O_j/\mathcal C_{O_j}$ is compact. Set
\[
 U=\bigcup_{j=1}^{m}\ \bigcup_{h\in H}hO_j,
 \qquad X=\overline U.                                  \tag{C.55.2}
\]
There is a based contraction on $X^\ddagger\cup X$,
jointly continuous with its inverse, commuting with $H$,
and equal to the prescribed contraction on $X^\ddagger$.
On every nonopen orbit in $X$ it has the exact hybrid
value of C.52.1. The intrinsic contraction on each
representative commutes with all of $\mathcal C_{O_j}$.
Thus these conclusions apply with the full local
centralizer quotient, without requiring that its
elements extend to the plane.
\end{claim}

\begin{proof}
\emph{The local group, its topology, and its global comparisons.}
The centralizer of $G|_O$ is also the centralizer of its
closure in $\operatorname{Homeo}(O)$. The latter is a
closed Polish group transitive on $O$. C.30 therefore
proves that $\mathcal C_O$ is locally compact, with a
free proper action. C.33's surface Lie argument applies
to this very action and makes $\mathcal C_O$ a Lie group
of dimension at most two. These conclusions use the
restriction image and its topology, not a discrete
topology on the original operators. Formula (C.55.1)
and its kernel follow directly from commutation on the
plane. Closure in $\operatorname{Homeo}(O)$ proves the
assertion about $Q_O$.

The global group $L$ preserves $O$, by the boundary
and inverse-limit argument of C.36. Its restriction is
continuous into $\operatorname{Homeo}(O)$: compact subsets
of $O$ and their inverse images remain compact subsets
of the plane. Hence every $q\in\mathcal C_O$ satisfies
\[
                      q(lu)=lq(u)\quad(l\in L,\ u\in O).
                                                        \tag{C.55.3}
\]
For example, choose $g_n\in G$ converging to $l$ globally,
restrict on compact subsets of $O$, and pass
$qg_nu=g_nqu$ to the limit. This proves (C.55.3)
without a plane extension of $q$.

For any $h\in H$, conjugation defines the topological
group isomorphism
\[
 \mathcal C_O\longrightarrow\mathcal C_{hO},\qquad
                         q\longmapsto hqh^{-1}.
                                                        \tag{C.55.4}
\]
The domain and codomain here are local homeomorphism
groups; the displayed conjugation uses only the actual
global operator $h$. It preserves dimension and the
compactness of the corresponding orbit quotient.

Fix $x\in O$, put $J=L_x$, and choose $a_q\in L$ with
$a_qx=qx$, using transitivity. Equation (C.55.3) and its
inverse give $L_{qx}=J$, and therefore
$a_qJa_q^{-1}=J$. For every global quotient in C.53.1,
\[
 \begin{split}
 \Phi_z^x q(lx)&=\Phi_z^x(la_qx)=la_qz,\\
                  \Phi_z^x q&=\Phi_{a_qz}^x.
 \end{split}                                            \tag{C.55.5}
\]
The target $a_qz$ is fixed by $J$ and lies in the same
actual nonopen orbit, since C.42 gives $Lz=Gz$.
Changing $a_q$ by its stabilizer ambiguity changes no
map in (C.55.5). Applying the formula also to $q^{-1}$
proves that the full global family is permuted by
$\mathcal C_O$. The exact transport between different
orbits is still C.53.5. All source and target coverings
in that family are therefore the ones already used
in C.43--C.48.

\emph{The enlarged exponential region.}
For $\dim\mathcal C_O=2$, the proof of C.36 applies with
$\mathcal C_O$ as its proper group. More explicitly, its
identity-component orbits are open by the free
dimension-two action and closed by properness; since
$O$ is connected, that component is transitive. Freeness
then gives $\mathcal C_O=\mathcal C_O^0$. The evaluation
map identifies it with $O$. Its full centralizer
\[
 N_O=C_{\operatorname{Homeo}(O)}(\mathcal C_O)
                            \cong\mathcal C_O^{\rm op}
                                                        \tag{C.55.6}
\]
is a simply transitive two-dimensional Lie group.
The actual restriction image $G|_O$ equals $N_O$:
transitivity supplies a word with any prescribed value
at one point, and evaluation is injective on $N_O$.

The resulting surjection $L\to N_O$ is continuous and
open. Its kernel fixes $\overline O$ pointwise, and it
fixes the image of every global quotient $\Phi_z$
pointwise by $l\Phi_z(u)=\Phi_z(lu)$. The two actions
therefore factor through the open quotient, with their
actual topology, exactly as in C.36.3 and C.53.6--C.53.7.
The universal covering group is $\mathbb R^2$ or
$\operatorname{Aff}^+(\mathbb R)$, with global exponential
and logarithm. The original based path lifts to a
closed path: C.32 gives its null-homotopy in
$\operatorname{Homeo}(O)$, and the evaluation
homeomorphism gives the identical statement in $N_O$.
Fractional powers in that specified cover define the
same closure contraction as C.36.6. On each global
nonopen quotient they give the exponential branch of
C.52.1, including the based real circle angle.

For $|\mathscr F_O^{\rm all}|\geq3$, retain exactly the
actual group and fractional powers of C.53.6--C.53.8.
They also commute with every $q\in\mathcal C_O$.
Indeed choose any lift $\widetilde q$ to the universal
cover of $O$. The path
$\widetilde q\widetilde f_t\widetilde q^{-1}$ is the
based lift of the same path $f_t|_O$, so uniqueness of
path lifting identifies it with $\widetilde f_t$.
It follows that $\widetilde q$ commutes with its entire
actual generated source group and with the displayed
fractional powers. Descent proves the assertion.
Where the two constructions apply together, they
use the same based path in the same actual simply
connected Lie group and agree.

All these source groups have the uniform fractional
word bound $1913$ of C.50; the two-dimensional groups
have the stronger bound three. Their images are
metabelian. The compact global word-isotopy family
C.50.15 consequently applies to their entire union.
For each positive integer $n$, its fractional residual
is in the same specified simply connected group and
has the same bound. On a limiting line orbit, the
bounded logarithm argument C.50.17--C.50.18 determines
the fractional limit. On a limiting circle orbit,
the null residual word loop has zero two-point winding,
so C.50.19 retains its real angle, including the
integer part. On fixed points every such word is the
identity. These are precisely the inputs of C.50's
dense-set limit and diagonal arguments. They prove
joint continuity on $X^\ddagger$ and, on applying them
to the inverse loop, joint inverse continuity.
Equation (C.55.4) and uniqueness of the based lift
give the required operator transport and commutation.

Since $Q_O$ is a closed subgroup of $\mathcal C_O$,
$\dim Q_O=2$ implies $\dim\mathcal C_O=2$. Thus
$U^\dagger\subset U^\ddagger$. The two contractions
agree on $U^\dagger$ by their actual source formulas,
and hence on $X^\dagger$ by continuity. This proves
the first extension assertion with its exact
prescription.

\emph{Intrinsic maps commuting with the full local group.}
Each remaining representative in (C.55.2) has
$\dim\mathcal C_{O_j}\leq1$ and at most two global
partitions. Apply the intrinsic construction of C.52
to the proper Lie group $\mathcal C_{O_j}$ and to this
full family. Every required operator property is now
proved: freeness, properness and the topology follow
from C.30; the Lie property follows from C.33; and
permutation of all quotient partitions is (C.55.5).
In the two-partition case, C.46 and C.52 use the exact
coordinate changes and target conjugacies. In the
one-partition case, C.48's proper-group normal forms
and C.52's fiber and base formulas retain every
relation of this very local group, including deck
relations. With an empty family, C.32--C.33 give the
intrinsic centralizer contraction using the same
ambient-null restriction. These intrinsic arguments
do not use a global extension of an operator.

They give a based family $K^j_{s,t}$ on $O_j$,
jointly continuous with its inverse, commuting with
$\mathcal C_{O_j}$, and satisfying
\[
              \Phi_z K^j_{s,t}=R^{Gz}_{s,t}\Phi_z
                  \quad\text{for every global quotient.}
                                                        \tag{C.55.7}
\]
Constant equivariant maps satisfy the identical
equation, with the fixed-point target prescription.
Remove duplicate $H$-representatives and set
$K=hK^jh^{-1}$ on $hO_j$. Its independence from the
chosen $h$ follows from $H_{O_j}|_{O_j}\subset
\mathcal C_{O_j}$. The union is a jointly continuous
homeomorphism family on the open set $U$, with
jointly continuous inverse and exact $H$-commutation.

\emph{Limits of the actual local equivariant maps.}
Compactness of $O_j/\mathcal C_{O_j}$ supplies a compact
$C_j\subset O_j$ with $\mathcal C_{O_j}C_j=O_j$.
For completeness, the quotient projection is open.
Relatively compact open neighborhoods in $O_j$ have
open quotient images covering the compact quotient;
a finite subcover and the union of their compact
closures give $C_j$.

Take $x_n\in U$ converging to $z\in X\setminus U$,
and $(s_n,t_n)\to(s,t)$. The orbit $Gz$ is nonopen,
because an open orbit containing $z$ would contain
all late $x_n$ and would be a component of $U$.
After a subsequence, the finite representative set
and compactness give one $j$ and
\[
 x_n=h_nq_ny_n,
 \quad h_n\in H,\ q_n\in\mathcal C_{O_j},\ y_n\in C_j,
 \qquad y_n\longrightarrow y\in C_j.
                                                        \tag{C.55.8}
\]
The maps
\[
 m_n:O_j\longrightarrow\mathbb R^2,
                  \qquad m_n(u)=h_nq_nu
                                                        \tag{C.55.9}
\]
are actual continuous $L$-equivariant maps by
(C.55.3) and global commutation of $h_n$ with $L$.
The open $L$-orbit map supplies $l_n\to1$ in $L$
with $l_ny=y_n$. Consequently
\[
                    m_n(y)=l_n^{-1}x_n\longrightarrow z.
                                                        \tag{C.55.10}
\]
Every element of $L_y$ fixes $z$. C.42's evaluation
homeomorphism, applied to the space of all continuous
$L$-equivariant maps $O_j\to\mathbb R^2$, gives
\[
                  m_n\longrightarrow\Phi_z^y
                         \quad\text{uniformly on compact sets}.
                                                        \tag{C.55.11}
\]
This is a direct use of the actual maps (C.55.9).
No element $q_n$ has been extended to the plane or
approximated by restrictions of global operators.

The points $K^j_{s_n,t_n}(y_n)$ converge to
$K^j_{s,t}(y)$ in a compact subset of $O_j$, by joint
continuity and compactness of the parameter square
times $C_j$. Transport, commutation, (C.55.11), the
basepoint comparison C.53.3, and (C.55.7) give exactly
\[
 \begin{split}
 K_{s_n,t_n}(x_n)
   &=h_nq_nK^j_{s_n,t_n}(y_n)\\
   &\longrightarrow\Phi_z^y K^j_{s,t}(y)
                      =R^{Gz}_{s,t}(z).
 \end{split}                                            \tag{C.55.12}
\]
Every subsequence has a further subsequence of this
form, and its limit is the same prescribed value.
Thus (C.55.12) holds for the whole sequence.
An empty global family makes each such limit orbit
a fixed point, since otherwise (C.55.11) itself
would supply a member of that family.

Define $K$ on $X\setminus U$ by this prescribed
hybrid value. At a varying boundary point and fixed
parameter, (C.55.12) supplies a nearby point of $U$
with nearby image. Choosing both distances less
than $1/n$ and applying (C.55.12) once more proves
joint continuity for boundary-only sequences.
Combining interior and boundary subsequences proves
joint continuity on the whole square times $X$.
The inverse form of (C.55.7) and the compact set of
all $(K^j_{s,t})^{-1}(C_j)$ give the identical argument
for inverses. Every map preserves the treated orbits
and has this displayed inverse; it is therefore a
homeomorphism of $X$. The based-edge identities and
$H$-commutation extend from the dense open set $U$.

Finally $X\cap X^\ddagger$ contains only nonopen
orbits. The metabelian identity holds on
$X^\ddagger$, so the hybrid value there is exactly
its prescribed exponential value, with the same
based covering lift. The two families and their
inverses agree on the intersection. Finite closed-set
pasting gives the contraction on $X^\ddagger\cup X$.

This argument still requires finitely many
$H$-representatives and compact quotients by their
full local centralizers. It supplies no extension
through arbitrary noncompact local quotients,
infinitely many inequivalent representatives, or
the remaining exterior nonopen region. Those parts
of C.12 remain under investigation.
\end{proof}
